\section{Problem Formulation: Joint Unknown-Space Exploration and Reconstruction}
\label{problem-formulation}

We consider a ground robot that must discover a previously unmapped static environment and reconstruct its observed surfaces. The robot starts from one acquired RGB-D and planar-scan packet and subsequently receives measurements only after executed actions. It is not supplied with the environment's complete occupancy map, a complete safe pose graph, hidden object bounds, or source meshes. The online state contains the current pose, an acquired occupancy map, a measured surface representation, acquired instance/category support, and an action/observation ledger. Reference geometry and the start-reachable free-space set are reserved for evaluation. Exact simulated poses and partial ideal first-hit instance/category labels are controlled sensing assumptions, not evidence of automatic detection or SLAM accuracy.

The primary objective is efficient unknown-space coverage together with measured surface reconstruction. Discovering new free space remains an objective throughout the exploration task; object completion supplies an additional observation preference rather than replacing exploration with inspection of a supplied target. Let $\mathcal D_t$ denote all measurements acquired by time $t$, $M_t$ the corresponding occupancy map, and $\mathcal Q_t$ the measured surface state. The declared v10\_3\_1-derived candidate generator produces a finite set $\mathcal R_t$ of observation routes from this acquired state. Its occupancy, connectivity, footprint, action-budget and return checks are shared by G and S. The two policies therefore compare the same routes at an identical acquired state. Once their actions diverge, their later observations and candidate sets may differ.

For a candidate $r$, the implemented joint score is
\begin{equation}
U_t^m(r)=J_t(r)+H_t^m(r),\qquad m\in\{G,S\},
\end{equation}
where $J_t$ retains the inherited unknown-coverage and measured-surface-quality terms and $H_t^m$ is the expected unobserved template area. More explicitly, the common coverage contribution is $a_{\rm cell}[\widehat\eta_{L,t}|\mathcal U_{L,t}(r)|+\widehat\eta_{D,t}|\mathcal U_{D,t}(r)|]$: the route masks $\mathcal U$ contain predicted new support for the planar scan and camera, and the inherited observation-gain update supplies $\widehat\eta$. The measured-quality proxy is added using the same coefficient and implementation in every matched arm. This is a total outbound-gain score; it does not divide the score by route length or replace it with the final reconstruction metric.

Completion is defined only for the shared acquired positive instance tracks. For object $o$, $h\in\{\text{closed},\text{open}\}$ denotes a structural hypothesis and $k$ denotes a geometric particle. Both policies use the same acquired support, conditional geometry, visibility model and observation history. G has a neutral structural-class prior, whereas S uses a finite prior from the acquired category votes. Thus G is a neutral class-prior control with the same instance frontend, not a completely semantics-free perception system. In the particle revision, the expected completion term integrates area-weighted visibility over the joint structure--geometry distribution. For each template point, the maximum over committed outbound successor poses forms a route union; previously measured or strictly free-space-contradicted points do not earn new area. Hypothesized surfaces never enter the measured TSDF. Geometry, likelihood-scale and structural-temperature revisions are identified separately in the method and matched experiments.

The planner selects the highest positive score, breaking ties by lower route cost and then candidate identifier. Its original heading-state local planner executes the selected option through forward and turn primitives. Each executed primitive is charged once, checked by the common execution layer, and followed by a real simulator observation and measured-map update. Completion of an option or its invalidation triggers another global selection; no future branch measurement is supplied to the decision. A trained reachability network is not active in these virtual experiments. The complete virtual comparisons use 0.2\,m forward moves, $90^{\circ}$ turns and a common 256-action budget that includes return to the initial pose. These are declared evaluation conditions, not a claim that the current physical deployment uses the same budget or return contract.

We evaluate coverage efficiency separately from reconstruction quality. With the evaluation-only start-reachable free set $\mathcal F_{\rm ref}$,
\begin{equation}
C_t^{\rm free}=\frac{|\{c\in\mathcal F_{\rm ref}:M_t(c)=\text{free}\}|}{|\mathcal F_{\rm ref}|},\qquad
F_{1,\tau}=\frac{2P_\tau R_\tau}{P_\tau+R_\tau}.
\end{equation}
The primary coverage outcomes are the first paid actions reaching declared thresholds and the complete coverage curve. We also report travel, action count, collisions, return outcome, correctly free and falsely occupied cells, whole-scene surface precision/recall/F1, and object recall. Surface measurements compare the native measured TSDF with a separately fixed visible source reference at $\tau=0.05$\,m (with the additional threshold retained in the supplement), without aligning or inserting predicted geometry. Object-surface improvement, whole-scene improvement and coverage improvement are distinct outcomes. Neither a changed category probability nor a changed target alone establishes an efficiency benefit.

